% One-time imported draft. Editable native source; no Markdown regeneration.
\section{Step 17 Verify the protocol and explain the
  bounds}\label{005-seven-segment-counter-s17-step-17-verify-the-protocol-and-explain-the-bounds}

The grader observes real execution writes. It checks every transaction,
including repeated values, and observes protected data for writes that a
final-value check could miss. This evidence is stronger than a
student-maintained RAM log of intended output.

\subsection{Write or
  inspect}\label{005-seven-segment-counter-s17-write-or-inspect}

Run the public checks. Explain the loop invariant, the final pointer,
the difference between an index and a mask, and why a single final-value
write is insufficient.

\subsection{Expected
  observation}\label{005-seven-segment-counter-s17-expected-observation}

The required sequence, port, width, and protected-data behavior must all
agree. UI rendering and repaint readiness remain separate integration
and capture checks.

\subsection{Check your
  understanding}\label{005-seven-segment-counter-s17-check-your-understanding}

Why is a RAM log written by the student program insufficient by itself?

\begin{itemize}
  \tightlist
  \item
        \textbf{A} RAM cannot hold byte values.
  \item
        \textbf{B} Every data store automatically drives the right display.
  \item
        \textbf{C} It can describe intended output without proving which
        device transactions occurred.
\end{itemize}
